\subsubsection{Auxiliary Payloads}
For an AUXILIARY frame, (:ref:base.events.frame.EVENT_FRAME.ERROR_CODE:) bit 26 is the applicable
(:ref:base.events.payloads.DELIVERY_FAILURE_ERROR_CODE.EVENT_AUX_VALID:),
(:ref:base.events.payloads.MACHINE_CHECK_ERROR_CODE.EVENT_AUX_VALID:), or
(:ref:base.events.payloads.BUS_FAILURE_ERROR_CODE.EVENT_AUX_VALID:) field. Bits 24 and 25 are the corresponding address-validity fields when the event can supply those addresses. Every
unassigned bit is zero, and an auxiliary or address value not produced is zero with its validity bit clear.

Delivery-failure stages are represented by distinct leaf events. Their (:ref:base.events.frame.EVENT_FRAME.EVENT_INFO.EVENT_CODE:) identifies the failed stage; the
(:ref:base.events.frame.EVENT_FRAME.EVENT_AUX:) (:term:base.terms.payload:) records the event whose delivery failed.

\begin{BedrockListedFormatDiagram}{Delivery-Failure Auxiliary Format}
\BedrockFormatRowRange{ERROR\_CODE[63:0]}{63}{0}{%
\BedrockFormatReserved{(:term:base.terms.reserved:)}{37}
\BedrockFormatField{E}{1}
\BedrockFormatField{L}{1}
\BedrockFormatField{F}{1}
\BedrockFormatReserved{(:term:base.terms.reserved:)}{24}
}
\BedrockFormatRowRange{EVENT\_AUX[63:0]}{63}{0}{%
\BedrockFormatField{zero}{32}
\BedrockFormatField{failed event code}{32}
}
\end{BedrockListedFormatDiagram}
{\small
\texttt{E}, \texttt{L}, and \texttt{F} are (:event-structure-target:base.events.payloads.DELIVERY_FAILURE_ERROR_CODE.EVENT_AUX_VALID:)\texttt{EVENT\_AUX\_VALID},
(:event-structure-target:base.events.payloads.DELIVERY_FAILURE_ERROR_CODE.FAULT_LINEAR_VALID:)\texttt{FAULT\_LINEAR\_VALID}, and (:event-structure-target:base.events.payloads.DELIVERY_FAILURE_ERROR_CODE.FAULT_EA_VALID:)\texttt{FAULT\_EA\_VALID}.
\par}

When (:ref:base.events.payloads.DELIVERY_FAILURE_ERROR_CODE.EVENT_AUX_VALID:) is one, (:ref:base.events.frame.EVENT_FRAME.EVENT_AUX:)\texttt{[31:0]} contains the event code whose delivery failed and bits 63..32 are zero. A
(:ref:base.events.EXCEPTION.EVENT_FRAME_ADDRESS_FAILURE:) or (:ref:base.events.EXCEPTION.EVENT_FRAME_STORE_FAILURE:) supplies (:ref:base.events.frame.EVENT_FRAME.FAULT_EA:) and (:ref:base.events.frame.EVENT_FRAME.FAULT_LINEAR:)
when those values were produced.

For (:ref:base.events.EXCEPTION.MACHINE_CHECK:), the error code has this layout:

\begin{BedrockListedFormatDiagram}{MACHINE\_CHECK ERROR\_CODE Format}
\BedrockFormatRowRange{ERROR\_CODE[63:0]}{63}{0}{%
\BedrockFormatReserved{(:term:base.terms.reserved:)}{37}
\BedrockFormatField{E}{1}
\BedrockFormatField{L}{1}
\BedrockFormatField{F}{1}
\BedrockFormatReserved{(:term:base.terms.reserved:)}{12}
\BedrockFormatField{R}{1}
\BedrockFormatField{P}{1}
\BedrockFormatField{SV}{2}
\BedrockFormatField{SRC}{8}
}
\end{BedrockListedFormatDiagram}
{\small
\texttt{E}, \texttt{L}, and \texttt{F} are (:ref:base.events.payloads.MACHINE_CHECK_ERROR_CODE.EVENT_AUX_VALID:),
(:ref:base.events.payloads.MACHINE_CHECK_ERROR_CODE.FAULT_LINEAR_VALID:), and (:ref:base.events.payloads.MACHINE_CHECK_ERROR_CODE.FAULT_EA_VALID:);
\texttt{P}, \texttt{R}, \texttt{SV}, and \texttt{SRC} denote (:ref:base.events.payloads.MACHINE_CHECK_ERROR_CODE.PRECISE:),
(:ref:base.events.payloads.MACHINE_CHECK_ERROR_CODE.RETRY_SAFE:), (:ref:base.events.payloads.MACHINE_CHECK_ERROR_CODE.SEVERITY:), and (:ref:base.events.payloads.MACHINE_CHECK_ERROR_CODE.SOURCE:).
\par}

\BedrockTableCaption{MACHINE\_CHECK ERROR\_CODE Fields}
\begin{BedrockTabular}{@{}>{\raggedright\arraybackslash}p{0.70in}>{\raggedright\arraybackslash}p{1.35in}>{\raggedright\arraybackslash}p{3.45in}@{}}
\toprule
\rowcolor{BedrockHeaderFill}
\textbf{Bits} & \textbf{Field} & \textbf{Meaning}\\
\midrule
7..0 & (:event-structure-target:base.events.payloads.MACHINE_CHECK_ERROR_CODE.SOURCE:)\texttt{SOURCE} & 0 unknown; 1 core; 2 cache; 3 translation; 4 memory; 5 interconnect\\
9..8 & (:event-structure-target:base.events.payloads.MACHINE_CHECK_ERROR_CODE.SEVERITY:)\texttt{SEVERITY} & 0 corrected; 1 recoverable; 2 fatal; 3 (:term:base.terms.reserved:)\\
10 & (:event-structure-target:base.events.payloads.MACHINE_CHECK_ERROR_CODE.PRECISE:)\texttt{PRECISE} & saved PC identifies the responsible instruction and saved (:term:base.terms.architectural_state:) precedes it\\
11 & (:event-structure-target:base.events.payloads.MACHINE_CHECK_ERROR_CODE.RETRY_SAFE:)\texttt{RETRY\_SAFE} & retry at that precise boundary cannot duplicate an external effect\\
23..12 & (:term:base.terms.reserved:) & zero\\
24 & (:event-structure-target:base.events.payloads.MACHINE_CHECK_ERROR_CODE.FAULT_EA_VALID:)\texttt{FAULT\_EA\_VALID} & (:ref:base.events.frame.EVENT_FRAME.FAULT_EA:) was produced\\
25 & (:event-structure-target:base.events.payloads.MACHINE_CHECK_ERROR_CODE.FAULT_LINEAR_VALID:)\texttt{FAULT\_LINEAR\_VALID} & (:ref:base.events.frame.EVENT_FRAME.FAULT_LINEAR:) was produced\\
26 & (:event-structure-target:base.events.payloads.MACHINE_CHECK_ERROR_CODE.EVENT_AUX_VALID:)\texttt{EVENT\_AUX\_VALID} & (:ref:base.events.frame.EVENT_FRAME.EVENT_AUX:) was produced\\
63..27 & (:term:base.terms.reserved:) & zero\\
\bottomrule
\end{BedrockTabular}

When (:ref:base.events.payloads.MACHINE_CHECK_ERROR_CODE.EVENT_AUX_VALID:) is one, (:ref:base.events.frame.EVENT_FRAME.EVENT_AUX:) contains an implementation-defined syndrome. Associated effective and linear
addresses are supplied when available.

(:ref:base.events.payloads.MACHINE_CHECK_ERROR_CODE.RETRY_SAFE:)\texttt{=1} requires (:ref:base.events.payloads.MACHINE_CHECK_ERROR_CODE.PRECISE:)\texttt{=1}. The valid (:ref:base.events.payloads.MACHINE_CHECK_ERROR_CODE.SEVERITY:) and (:term:base.terms.continuation:) combinations are:

\BedrockTableCaption{MACHINE\_CHECK Valid Continuation Combinations}
\begin{BedrockLongTable}{@{}>{\raggedright\arraybackslash}p{0.95in}>{\raggedright\arraybackslash}p{0.55in}>{\raggedright\arraybackslash}p{0.75in}>{\raggedright\arraybackslash}p{3.25in}@{}}
\toprule
\rowcolor{BedrockHeaderFill}
\textbf{(:ref:base.events.payloads.MACHINE_CHECK_ERROR_CODE.SEVERITY:)} & \textbf{(:ref:base.events.payloads.MACHINE_CHECK_ERROR_CODE.PRECISE:)} & \textbf{(:ref:base.events.payloads.MACHINE_CHECK_ERROR_CODE.RETRY_SAFE:)} & \textbf{Saved PC and (:term:base.terms.continuation:) meaning}\\
\midrule
\endfirsthead
\multicolumn{4}{l}{\scriptsize\itshape Table \theBedrockTable\ (continued)}\\
\toprule
\rowcolor{BedrockHeaderFill}
\textbf{(:ref:base.events.payloads.MACHINE_CHECK_ERROR_CODE.SEVERITY:)} & \textbf{(:ref:base.events.payloads.MACHINE_CHECK_ERROR_CODE.PRECISE:)} & \textbf{(:ref:base.events.payloads.MACHINE_CHECK_ERROR_CODE.RETRY_SAFE:)} & \textbf{Saved PC and (:term:base.terms.continuation:) meaning}\\
\midrule
\endhead
corrected & 0 & 0 & Trap; saved PC is the instruction following the corrected operation.\\
recoverable & 0 & 0 & Saved PC is the next instruction that would execute at the delivery boundary; saved (:term:base.terms.architectural_state:) is the state at that boundary. It does not identify the responsible operation.\\
recoverable & 1 & 0 & Saved PC identifies the responsible instruction and (:term:base.terms.architectural_state:) precedes it; retry may duplicate an external effect.\\
recoverable & 1 & 1 & Saved PC identifies the responsible instruction and (:term:base.terms.architectural_state:) precedes it; retry cannot duplicate an external effect.\\
fatal & 0 & 0 & Saved PC is the next instruction at the delivery boundary and is diagnostic only.\\
fatal & 1 & 0 & Saved PC identifies the responsible instruction and pre-instruction (:term:base.terms.architectural_state:) for diagnosis only.\\
\bottomrule
\end{BedrockLongTable}

All other combinations are invalid and are not generated. In particular, corrected machine checks never set either
bit, (:ref:base.events.payloads.MACHINE_CHECK_ERROR_CODE.PRECISE:)\texttt{=0}, (:ref:base.events.payloads.MACHINE_CHECK_ERROR_CODE.RETRY_SAFE:)\texttt{=1} is invalid for every severity, fatal machine checks never set
(:ref:base.events.payloads.MACHINE_CHECK_ERROR_CODE.RETRY_SAFE:), and (:ref:base.events.payloads.MACHINE_CHECK_ERROR_CODE.SEVERITY:) 3 is (:term:base.terms.reserved:). A fatal event provides no reliable (:term:base.terms.continuation:) even when its
diagnostic state is precise.

For an event in the \texttt{BUS\_FAILURE} family, the error code has this layout:

\begin{BedrockListedFormatDiagram}{BUS\_FAILURE ERROR\_CODE Format}
\BedrockFormatRowRange{ERROR\_CODE[63:0]}{63}{0}{%
\BedrockFormatReserved{(:term:base.terms.reserved:)}{33}
\BedrockFormatField{R}{1}
\BedrockFormatField{SZ}{3}
\BedrockFormatField{E}{1}
\BedrockFormatField{L}{1}
\BedrockFormatField{F}{1}
\BedrockFormatField{OP}{8}
\BedrockFormatReserved{0}{1}
\BedrockFormatField{LVL}{3}
\BedrockFormatField{A}{1}
\BedrockFormatField{U}{1}
\BedrockFormatField{AC}{2}
\BedrockFormatReserved{(:term:base.terms.reserved:)}{8}
}
\end{BedrockListedFormatDiagram}
{\small
\texttt{E}, \texttt{L}, and \texttt{F} are (:ref:base.events.payloads.BUS_FAILURE_ERROR_CODE.EVENT_AUX_VALID:),
(:ref:base.events.payloads.BUS_FAILURE_ERROR_CODE.FAULT_LINEAR_VALID:), and (:ref:base.events.payloads.BUS_FAILURE_ERROR_CODE.FAULT_EA_VALID:);
\texttt{SZ}, \texttt{OP}, \texttt{LVL}, and \texttt{AC} denote (:ref:base.events.payloads.BUS_FAILURE_ERROR_CODE.ACCESS_SIZE:),
(:ref:base.events.payloads.BUS_FAILURE_ERROR_CODE.OPERAND:), (:ref:base.events.payloads.BUS_FAILURE_ERROR_CODE.WALK_LEVEL:), and (:ref:base.events.payloads.BUS_FAILURE_ERROR_CODE.ACCESS:); \texttt{A}, \texttt{U}, and
\texttt{R} denote (:ref:base.events.payloads.BUS_FAILURE_ERROR_CODE.ATOMIC:), (:ref:base.events.payloads.BUS_FAILURE_ERROR_CODE.USER_DOMAIN:), and (:ref:base.events.payloads.BUS_FAILURE_ERROR_CODE.RETRY_SAFE:).
\par}

\BedrockTableCaption{BUS\_FAILURE ERROR\_CODE Fields}
\begin{BedrockLongTable}{@{}>{\raggedright\arraybackslash}p{0.70in}>{\raggedright\arraybackslash}p{1.35in}>{\raggedright\arraybackslash}p{3.45in}@{}}
\toprule
\rowcolor{BedrockHeaderFill}
\textbf{Bits} & \textbf{Field} & \textbf{Meaning}\\
\midrule
\endfirsthead
\multicolumn{3}{l}{\scriptsize\itshape Table \theBedrockTable\ (continued)}\\
\toprule
\rowcolor{BedrockHeaderFill}
\textbf{Bits} & \textbf{Field} & \textbf{Meaning}\\
\midrule
\endhead
7..0 & (:term:base.terms.reserved:) & zero; failure identity is carried by (:ref:base.events.frame.EVENT_FRAME.EVENT_INFO.EVENT_CODE:)\\
9..8 & (:event-structure-target:base.events.payloads.BUS_FAILURE_ERROR_CODE.ACCESS:)\texttt{ACCESS} & 0 none, 1 read, 2 write, 3 execute\\
10 & (:event-structure-target:base.events.payloads.BUS_FAILURE_ERROR_CODE.USER_DOMAIN:)\texttt{USER\_DOMAIN} & access selected the user-domain permission path\\
11 & (:event-structure-target:base.events.payloads.BUS_FAILURE_ERROR_CODE.ATOMIC:)\texttt{ATOMIC} & access was an atomic read-modify-write\\
14..12 & (:event-structure-target:base.events.payloads.BUS_FAILURE_ERROR_CODE.WALK_LEVEL:)\texttt{WALK\_LEVEL} & 0 outside a walk; 1 through 4 identify the PTE level\\
15 & (:term:base.terms.reserved:) & zero\\
23..16 & (:event-structure-target:base.events.payloads.BUS_FAILURE_ERROR_CODE.OPERAND:)\texttt{OPERAND} & zero-based explicit operand ordinal; \texttt{0xff} for implicit fetch or stack access\\
24 & (:event-structure-target:base.events.payloads.BUS_FAILURE_ERROR_CODE.FAULT_EA_VALID:)\texttt{FAULT\_EA\_VALID} & (:ref:base.events.frame.EVENT_FRAME.FAULT_EA:) was produced\\
25 & (:event-structure-target:base.events.payloads.BUS_FAILURE_ERROR_CODE.FAULT_LINEAR_VALID:)\texttt{FAULT\_LINEAR\_VALID} & (:ref:base.events.frame.EVENT_FRAME.FAULT_LINEAR:) was produced\\
26 & (:event-structure-target:base.events.payloads.BUS_FAILURE_ERROR_CODE.EVENT_AUX_VALID:)\texttt{EVENT\_AUX\_VALID} & (:ref:base.events.frame.EVENT_FRAME.EVENT_AUX:) was produced\\
29..27 & (:event-structure-target:base.events.payloads.BUS_FAILURE_ERROR_CODE.ACCESS_SIZE:)\texttt{ACCESS\_SIZE} & 0 unavailable or inapplicable; 1 B; 2 W; 3 L; 4 Q\\
30 & (:event-structure-target:base.events.payloads.BUS_FAILURE_ERROR_CODE.RETRY_SAFE:)\texttt{RETRY\_SAFE} & retry cannot duplicate an external effect\\
63..31 & (:term:base.terms.reserved:) & zero\\
\bottomrule
\end{BedrockLongTable}

A bus error during a page walk reports the level whose PTE access failed. When (:ref:base.events.payloads.BUS_FAILURE_ERROR_CODE.EVENT_AUX_VALID:) is one, (:ref:base.events.frame.EVENT_FRAME.EVENT_AUX:)
contains the final byte address available for the failed access.
